% Generated draft: requires mathematical and source review.
\documentclass[11pt]{article}
\usepackage[margin=25mm]{geometry}
\usepackage{fontspec}
\setmainfont{Latin Modern Roman}
\usepackage{amsmath,amssymb,mathtools,unicode-math}
\setmathfont{Latin Modern Math}
\usepackage{enumitem,needspace,xurl,hyperref}
\setlength{\parindent}{0pt}
\setlength{\parskip}{5pt}
\setlength{\emergencystretch}{3em}
\newcommand{\N}{\mathbb{N}}
\newcommand{\Z}{\mathbb{Z}}
\newcommand{\Q}{\mathbb{Q}}
\newcommand{\R}{\mathbb{R}}
\newcommand{\C}{\mathbb{C}}
\newcommand{\F}{\mathbb{F}}
\newcommand{\A}{\mathbb{A}}
\newcommand{\PP}{\mathbb{P}}
\newcommand{\E}{\mathbb{E}}
\newcommand{\eps}{\varepsilon}
\newcommand{\dd}{\,\mathrm{d}}
\begin{document}
\section*{UnsolvedMath ID: 30003954}
\textbf{Source problem code:} catalogue entry
\subsection{Definitions and mathematical statement}
Let \(S\) be a surface and let \(G=\operatorname{Mod}(S)\) denote its mapping class group: the group of isotopy classes of the permitted self-homeomorphisms of \(S\). The surface, orientation requirements, and conventions concerning boundary components and marked points must be specified before this defines a particular group; the supplied record does not specify them.

Let \(\mathcal C\) be the intended family of congruence subgroups of \(G\). Once \(S\), the mapping class group conventions, and \(\mathcal C\) are fixed, the source asks for a classification of the members of
\[
\mathcal N_1(G,\mathcal C)=\left\{N\trianglelefteq G: N\ne G,\quad (\forall C\in\mathcal C)\;N\nsubseteq C\right\},
\]
or, alternatively, of
\[
\mathcal N_2(G,\mathcal C)=\left\{N\trianglelefteq G: N\ne G,\quad N\nsubseteq\bigcup_{C\in\mathcal C}C\right\}.
\]
Here \(N\trianglelefteq G\) means that \(N\) is a normal subgroup of \(G\). The first condition means
\[
(\forall C\in\mathcal C)(\exists n\in N)\;n\notin C,
\]
whereas the alternative means
\[
(\exists n\in N)(\forall C\in\mathcal C)\;n\notin C.
\]
These quantifier orders express distinct requests and are retained separately.

The term ``congruence subgroup'' is not defined in the supplied statement. A possible homological level interpretation uses
\[
G[m]=\ker\!\left(G\longrightarrow\operatorname{Aut}\bigl(H_1(S;\Z/m\Z)\bigr)\right),\qquad m\ge2,
\]
provided the chosen mapping class group acts on this homology group, and takes \(\mathcal C\) to consist of these level kernels. This is only a conditional interpretation, not a specification established by the supplied source. The choice between level kernels and a broader family of congruence subgroups must be resolved. In particular, whether \(G\) itself is excluded from \(\mathcal C\) is essential: if \(G\in\mathcal C\), both displayed classes are empty.

The supplied record labels the problem partially solved and reports that Clay--Mangahas--Margalit construct explicit proper normal subgroups outside every level congruence subgroup, settling existence while leaving the requested classification open. This is a supplied dated assessment, not an independently verified mathematical or literature conclusion.
\subsection{Short English statement}
Classify the proper normal subgroups of a mapping class group that are contained in no congruence subgroup, or alternatively those that are not contained in the union of the congruence subgroups. The surface and the family of congruence subgroups are unspecified, and the two alternatives have different quantifier orders. The supplied partially solved label requires independent review.
\subsection{Sources}
\begin{itemize}
\item Supplied local UnsolvedMath registry record, permanent ID \(30003954\), title ``Noncongruence Normal Subgroups of Mapping Class Groups,'' exact source problem code \texttt{catalogue entry}, in the collection ``Oberwolfach Reports Open Problems.'' The record gives proposed year 2018 and no named proposer.
\item Supplied report citation: \emph{New Trends in Teichmüller Theory and Mapping Class Groups} (2019), Oberwolfach Reports, pp. 2475--2534, DOI \texttt{10.4171/owr/2018/40}. \url{https://doi.org/10.4171/owr/2018/40}. The background describes Oberwolfach Reports as documenting mathematical workshops and their open-problem discussions.
\item Supplied literature reference labeled ``Oberwolfach Report 40, Lanier problem'': \url{https://doi.org/10.4171/owr/2018/40}. This label supplies an association with Lanier but does not fill the record's unspecified proposer field.
\item Supplied literature reference labeled ``Clay--Mangahas--Margalit, RAAG normal subgroups'': \url{https://arxiv.org/abs/2001.10587}.
\item Literature assessment dated exactly \texttt{2026-08-21}: the supplied curation labels the problem partially solved and attributes explicit proper normal subgroups outside every level congruence subgroup to Clay--Mangahas--Margalit, while stating that classification remains open. This is supplied dated curation derived from the OWR material, not independently verified; its time-sensitive status and cited claims require independent verification.
\end{itemize}
\end{document}
